\scanpage{210}
SYSTÈMES de RACINES \mspage{(2)}\hfill Tribu n\textsuperscript{o} 59

Pour orienter le lecteur, nous indiquons ici les principaux exemples de systèmes de racines que l'on rencontre dans la théorie des groupes \unsure{$\cdot$}et algèbres de Lie.

I. Soient $\underline{\mathrm g}$ une algèbre de Lie semi-simple sur $\underline{\underline{\mathrm C}}$ et $\underline{\mathrm h}$ une sous-algèbre de Cartan de $\underline{\mathrm g}$. L'algèbre $\underline{\mathrm h}$ est donc commutative, formée d'éléments semi-simples. Une racine de $\underline{\mathrm g}$ par rapport à $\underline{\mathrm h}$ est une forme linéaire $\alpha$ non nulle sur $\underline{\mathrm h}$ pour laquelle il existe $\mathrm x\in\underline{\mathrm g}$, $\mathrm x\ne0$, tel que $[\mathrm h,\mathrm x]=\alpha(\mathrm x)$ quel que soit $\mathrm h\in\underline{\mathrm h}$. Soit $\mathrm V$ l'espace vectoriel sur $\underline{\underline{\mathrm Q}}$ engendré par l'ensemble $\mathrm R$ des racines de $\underline{\mathrm g}$ par rapport à $\underline{\mathrm h}$. Alors $\underline{\mathrm h}^{*}=\mathrm V\otimes\underline{\underline{\mathrm C}}$, et $\mathrm R$ est un système de racines réduit dans $\mathrm V$, au sens du \S~6. Tout système de racines réduit est l'ensemble des racines d'une algèbre semi-simple complexe, unique à son isomorphisme près. Tout automorphisme de $\underline{\mathrm g}$ laissant stable $\underline{\mathrm h}$ induit un automorphisme de $\underline{\mathrm h}^{*}$ laissant $\mathrm R$ et $\mathrm V$ invariants, et le groupe des automorphismes de $\mathrm R$ ainsi obtenus à partir des automorphismes, (resp. des automorphismes intérieurs), de $\underline{\mathrm g}$ s'identifie au groupe de tous les automorphismes de $\mathrm R$, (resp. au groupe de Weyl de $\mathrm R$)\unsure{:|.} Identifions le dual de $\mathrm V$ à une $\underline{\underline{\mathrm Q}}$-forme $\underline{\mathrm h}^{\unsure{\mathrm o|0|\circ}}$ de $\underline{\mathrm h}$ et soit $\mathrm G$ un groupe de Lie complexe connexe d'algèbres de Lie $\underline{\mathrm g}$. Alors le noyau $\Gamma$ de l'application exponentielle de $\underline{\mathrm h}$ dans $\mathrm G$ est un réseau de $\mathrm i.\underline{\mathrm h}^{\mathrm o}$ compris entre $2\pi\mathrm i.\mathrm Q(\mathrm R^{*})$ et $2\pi\mathrm i.\mathrm P(\mathrm R^{*})$, où $\mathrm R^{*}$ est le système de racines inverse de $\mathrm R$, qui est égal à $2\pi\mathrm i.\mathrm Q(\mathrm R^{*})$, (resp. $2\pi\mathrm i\mathrm P(\mathrm R^{*})$), si $\mathrm G$ est simplement connexe, (resp. isomorphe à $\mathrm{Ad}\,\underline{\mathrm g}$). Le groupe $2\pi\mathrm i.\mathrm P(\mathrm R^{*})/\Gamma$ s'identifie au centre de $\mathrm G$, et $\Gamma/2\pi\mathrm i.\mathrm Q(\mathrm R^{*})$ au groupe fondamental de $\mathrm G$. Le réseau de $\mathrm V$ associé à $(2\pi\mathrm i)^{-1}.\Gamma$ est engendré par les poids des représentations de $\mathrm G$.

II. Soient $\underline{\mathrm g}$ une algèbre de Lie semi-simple sur $\underline{\underline{\mathrm R}}$, $\underline{\mathrm k}$ une sous-algèbre compacte maximale de $\underline{\mathrm g}$, $\underline{\mathrm p}$ le complément orthogonal de $\underline{\mathrm k}$ par rapport à la forme de Killing, et $\underline{\mathrm a}$ une sous-algèbre commutative maximale de $\underline{\mathrm p}$. Pour toute forme
% archive sheet number: 209

